% Part: proof-theory
% Chapter: sequent-calculus
% Section: interpretation-rules

\documentclass[../../../include/open-logic-section]{subfiles}

\begin{document}

\olfileid{pt}{seq}{irl}

\olsection{నియమాల అర్థవ్యాఖ్యానం}

సీక్వెంట్ అర్థవ్యాఖ్యానాన్ని గుర్తుచేసుకోండి: $\Gamma \Sequent
\Delta$ అనేది \tetoken{ఒక నిర్మాణం}{!!a{structure}}~$\Struct{M}$లో సత్యమవ్వాలంటే, కనీసం ఒక $!A \in \Gamma$ అసత్యం కావాలి లేదా కనీసం ఒక~$!A \in \Delta$ సత్యం కావాలి. \Log{G3c} తార్కిక నియమాలను వాటి ప్రధాన \tetoken{సూత్రాల}{!!{formula}s} సత్య షరతులను సంకేతీకరించేవిగా చదవవచ్చు. \RightR{\lif}ను పరిగణించండి. దాని ప్రధాన \tetoken{సూత్రం}{!!{formula}} $!A \lif !B$; $!A$ అసత్యమైనా, $!B$ సత్యమైనా అది సత్యం. కాబట్టి $\ \Sequent !A \lif !B$ సత్యమైతేనే $!A \Sequent !B$ సత్యం. ఈ సీక్వెంట్ల పూర్వపక్షాలు, ఉత్తరపక్షాలకు సందర్భ \tetoken{సూత్రాలు}{!!{formula}s} $\Gamma$, $\Delta$లను చేరిస్తే, సరిగ్గా \RightR{\lif} నిర్ధారణ, ఆధారం వస్తాయి. \LeftR{\lif}కూ ఇలాంటిదే: $!A$ సత్యమై, $!B$ అసత్యమైతే $!A
\lif !B$ అసత్యం. కాబట్టి $!A \lif !B
\Sequent \ $ సత్యమైతేనే $\ \Sequent !A$, $!B \Sequent \ $ రెండూ సత్యం. \LeftR{\lif} నిర్ధారణ దాని ఆధారాల సంయోగానికి సమానం. ఈ పద్ధతి \Log{G3c} నియమాల సమర్థతను నిర్ధారిస్తుంది (ఆధారాలన్నీ సత్యమైతే నిర్ధారణ సత్యం); పరిపూర్ణతనూ నిర్ధారిస్తుంది.

నిజానికి ఏ సత్య-ప్రమేయ సంధానకానికైనా సత్య షరతులను ఇలా సీక్వెంట్ల సమితిగా వర్ణించవచ్చు. సాంప్రదాయ తర్కంలోని ప్రతి సత్య ప్రమేయాన్ని సంయుక్త సాధారణ రూపంలోని \tetoken{ఒక సూత్రంతో}{!!a{formula}} వర్ణించవచ్చనే వాస్తవం నుంచి ఇది అనుసరిస్తుంది: పరమాణు, నిషేధిత పరమాణు \tetoken{సూత్రాల}{!!{formula}s} వియోగాల సంయోగం. ఉదాహరణకు $!A
\oplus !B$ అనేది $!A$, $!B$లలో సరిగ్గా ఒక్కటి సత్యమైనప్పుడే సత్యమని తీసుకుందాం; అంటే “విశిష్ట వియోగం”. $(!A \lor !B) \land
(\lnot !A \lor \lnot !B)$ సత్యమైతేనే $!A \oplus !B$ సత్యం; $(\lnot !A \lor !B)
\land (!A \lor \lnot !B)$ సత్యమైతే అది అసత్యం. కాబట్టి $\oplus$కు సీక్వెంట్ నియమాలను ఇలా ప్రవేశపెట్టవచ్చు:
\begin{align*}
&
\Axiom$\Gamma \fCenter \Delta, !A, !B$
\Axiom$!A, !B, \Gamma \fCenter \Delta$
\RightLabel{\RightR{\oplus}}
\BinaryInf$\Gamma \fCenter \Delta, !A \oplus !B$
\DisplayProof
\\
& \Axiom$!A, \Gamma \fCenter \Delta, !B$
\Axiom$!B, \Gamma \fCenter \Delta, !A$
\RightLabel{\LeftR{\oplus}}
\BinaryInf$!A \oplus !B, \Gamma \fCenter \Delta$
\DisplayProof
\end{align*}
ఈ నియమాలు కూడా సమర్థమైనవి, పరిపూర్ణమైనవి. ఆధారాలన్నీ సత్యమైనప్పుడే వాటి నిర్ధారణలు సత్యం.
\sourcecorrection{OLTEPTSEQIRL-002}{రెండవ విశిష్ట వియోగ నియమం ప్రధాన సూత్రాన్ని ఎడమవైపు ఉంచినా మూలం దానికి కుడి నియమం గుర్తు ఇచ్చింది; ఎడమ నియమంగా గుర్తును సరిచేశాం.}

\begin{prob}
$!A$, $!B$ రెండూ సత్యం కానప్పుడే $!A \downarrow !B$ సత్యమని (“NOR”) అనుకుందాం. దాని \Log{G3c} నియమాలు ఎలా ఉంటాయి? ($!A \downarrow !B$ సత్యమైనప్పుడే ఒకేసారి సత్యమయ్యే రెండు సీక్వెంట్లను కనుగొనండి; అవి \RightR{\downarrow} నియమాన్ని ఇస్తాయి. తరువాత $!A \downarrow !B$ అసత్యమైనప్పుడే సత్యమయ్యే ఒక సీక్వెంట్‌ను కనుగొనండి; అది \LeftR{\downarrow} నియమాన్ని ఇస్తుంది.)
\end{prob}

\Log{G3c}లో ప్రతి సంధానకానికి, పరిమాణకానికి సరిగ్గా రెండు నియమాలు ఉంటాయి: ఎడమ నియమం, కుడి నియమం. కానీ \Log{G1c}లో రెండు \LeftR{\land}, రెండు \RightR{\lor} నియమాలు ఉన్నాయి. కారణం: \Log{G1c}కు నమూనా అయిన గెంట్సెన్ అసలు వ్యవస్థ~\Log{LK}కు, ముఖ్యమైన అసాంప్రదాయ తర్కం—అంతఃప్రజ్ఞావాద తర్కం—కోసం \Log{LJ} అనే రూపం ఉంది. అంతఃప్రజ్ఞావాద తర్కానికి సమర్థమైన, పరిపూర్ణమైన వ్యవస్థను పొందడానికి \Log{LJ}లో సీక్వెంట్ ఉత్తరపక్షంలో గరిష్ఠంగా ఒక \tetoken{సూత్రం}{!!{formula}} మాత్రమే అనుమతిస్తారు. అందువల్ల \Log{G3c} \RightR{\lor}ను వాడలేము; దాని ఆధార ఉత్తరపక్షంలో $!A$, $!B$ రెండూ ఉంటాయి. కాబట్టి గెంట్సెన్ \Log{LJ}లో రెండు \RightR{\lor} నియమాలను వాడి, \Log{LK}లోనూ అవే వాడారు. అలాగే అంతఃప్రజ్ఞావాద రూపం~\Log{G1i}, ఉత్తరపక్షాల్లో గరిష్ఠంగా ఒక \tetoken{సూత్రం}{!!{formula}} ఉండేలా పరిమితం చేసిన \Log{G1c}. (\Log{G3c}కూ అంతఃప్రజ్ఞావాద రూపం ఉంది; కానీ కొన్ని నియమాలు \Log{G3c}లోని అనురూప నియమాల నుంచి భిన్నంగా ఉంటాయి. ఉదాహరణకు అది కూడా రెండు \RightR{\lor} నియమాలను వాడుతుంది.)

\Log{G1c} నిర్మాణ నియమాలు సమర్థమైనవని సులభంగా చూడవచ్చు. ఉదాహరణకు $\Gamma \Sequent \Delta$లో కుడివైపు సత్యమైన \tetoken{సూత్రం}{!!{formula}} లేదా ఎడమవైపు అసత్యమైన \tetoken{సూత్రం}{!!{formula}} ఉంటే, అదే సూత్రం $\Gamma \Sequent
\Delta, !A$లోనూ ఉంటుంది; $!A$ సత్యమా అసత్యమా అన్నది సంబంధం లేదు. కాబట్టి \RightR{\Weakening} ఆధారం సత్యమైతే నిర్ధారణ కూడా సత్యం. ఇక్కడ వ్యతిరేక దిశ వర్తించదని గమనించండి: $!A$ సత్యమవడం వల్ల నిర్ధారణ సత్యం కావచ్చు; అప్పుడు ఆధారం సత్యమని హామీ లేదు.
\sourcecorrection{OLTEPTSEQIRL-001}{మూల బలహీనీకరణ వివరణ కుడివైపు అసత్యం, ఎడమవైపు సత్యం అని సీక్వెంట్ సత్య షరతులను తారుమారు చేసింది; పూర్వ నిర్వచనానికి అనుగుణంగా ధ్రువత్వాలను సరిచేశాం.}

అసలు నిర్మాణ నియమాలు ఎందుకు కావాలి? ఉదాహరణకు $\ \Sequent !A
\lor \lnot !A$ను \tetoken{నిరూపించడానికి}{!!{prove}} సంకోచన (\RightR{\Contraction}, \LeftR{\Contraction}) అవసరం. $!A \Sequent !A$తో మొదలుపెడితే, \RightR{\lnot}తో ముందుగా $\Sequent !A, \lnot !A$ వస్తుంది. తరువాత \RightR{\lor} రెండు రూపాలను వరుసగా వాడవచ్చు: ఒకసారి $\lnot !A$ నుంచి $!A \lor \lnot
!A$ను, మరొకసారి~$!A$ నుంచి $!A \lor \lnot !A$ను పొందుతాం. ఫలితం $\ \Sequent !A \lor \lnot !A, !A \lor \lnot !A$. కాబట్టి \Log{G1c}లో $\Sequent !A \lor \lnot !A$ను పొందడానికి, \tetoken{ఒక నిరూపణలో}{!!a{proof}} ఒకే \tetoken{సూత్రం}{!!{formula}} యొక్క రెండు ఉనికులను “సంకోచించగలగాలి”:
\[
\Axiom$!A \fCenter !A$
\RightLabel{\RightR{\lnot}}
\UnaryInf$\fCenter !A, \lnot !A$
\RightLabel{\RightR{\lor}}
\UnaryInf$\fCenter !A, !A \lor \lnot !A$
\RightLabel{\RightR{\lor}}
\UnaryInf$\fCenter !A \lor \lnot !A, !A \lor \lnot !A$
\RightLabel{\RightR{\Contraction}}
\UnaryInf$\fCenter !A \lor \lnot !A$
\DisplayProof
\]
నిజానికి బలహీనీకరణ లేదా సంకోచనం లేకుండా \Log{G1c}లో \tetoken{నిరూపించలేని}{!!{prove}} చెల్లుబాటు గల సీక్వెంట్లు ఉన్నాయి. ఉదాహరణకు $!A \Sequent !B \lif !A$కు \Log{G1c}లో \tetoken{ఒక నిరూపణను}{!!a{proof}} కింది విధంగా పొందుతాం. ఇక్కడ సూత్ర గుర్తులకు వేర్వేరు పరమాణు సూత్రాలను తీసుకుని, అదనపు సంకోచన దశలను పక్కనపెడితే, చివరి తార్కిక నియమం \RightR{\lif} కావాలి: ప్రధాన సూత్రం కాగలది $!B \lif !A$ మాత్రమే. దాని ఆధారం $!B, !A \Sequent !A$ కావాలి. దాన్ని అక్షయం $!A \Sequent !A$ నుంచి పొందడానికి \LeftR{\Weakening} అవసరం:
\[
\Axiom$!A \fCenter !A$
\RightLabel{\LeftR{\Weakening}}
\UnaryInf$!B, !A \fCenter !A$
\RightLabel{\RightR{\lif}}
\UnaryInf$!A \fCenter !B \lif !A$
\DisplayProof
\]
\sourcecorrection{OLTEPTSEQIRL-003}{మూలం ఈ నిరూపణ తప్పక షరతు కుడి నియమంతో ముగియాలని నిర్బంధం లేకుండా చెప్పింది. సూత్రాలు సంక్లిష్టమైనా, అదనపు నిర్మాణ దశలున్నా ఆ వాదన అక్షరాలా వర్తించదు. ఉద్దేశించిన వేర్వేరు పరమాణు సూత్రాల ఉదాహరణకు, చివరి తార్కిక నియమానికి వాదనను పరిమితం చేశాం.}

\Log{G3c}లో సంకోచన, బలహీనీకరణ నియమాలు అసలు అవసరం లేదు. బలహీనీకరణ అక్షయాల్లోనే అంతర్భాగం. రెండు ఆధారాల నియమాల్లో సందర్భం $\Gamma$, $\Delta$ రెండు ప్రతులను ఒకటిగా కలపడం వల్ల సంకోచనాన్ని చాలావరకు నివారించవచ్చు. \Log{G1c}, \Log{G3c} రెండింటిలోనూ ఇలాంటి “సందర్భాన్ని పంచుకునే” నియమాలు ఉన్నాయి. ఒక ఆధారం గల నియమాల్లో \RightR{\lor}, \LeftR{\land}, \RightR{\lexists}, \LeftR{\lforall}లకు మాత్రమే సంకోచనం అవసరం. \Log{G3c}లో ఆ నియమాల రూపాలు సంకోచనాన్ని పూర్తిగా అనవసరం చేస్తాయి. రెండు-ఆధార నియమాలకు స్వతంత్ర సందర్భాలు ఉండే \Log{G2c} వ్యవస్థ కూడా ఉంది (\olref{tab:G2c} చూడండి). \Log{G2c}లో సంకోచనం \Log{G1c} కంటే మరింత ముఖ్యమైనది.

\subfile{rules-G2c}

\end{document}
